% Pursuing Stacks (cotes 134-2 to 134-8): archive cote 134-5, site page 7 (PDF page 8), typescript page 194.
% typescript. Draft compared: ps4.tex:293-349.
% AI-transcribed by Claude Opus 5.5 (claude-opus-5-5, adaptive thinking, effort high), Message Batches API, batch ps2 (msgbatch_01NVwr1Yuo8qJf5rykeGBXPK), 2026-11-12; not yet checked by eye.
\scanpage{134-5/7}
\mspage{194}

I thought I was starting a retrospective of six weeks of scratchwork
on homotopical algebra  -  and it turns out to be a \inserted{(very)} short retrospective
rather of the up-and-down movement of \inserted{my} mathematical interests and invest-%
ments\hand{|}during the last years.  Doubtlessly the very strongest\unsure{\hand{,}} \inserted{attraction\unsure{,}} the greatest
fascination goes with this ``new world'' of anabelian algebraic geometry.
It may seem strange that instead,  I am indulging in this lengthy digressi-%
-on on homotopical algebra,  which is almost wholly irrelevant I feel for
the Galois-Teichm\xout{\unsure{u}}üller story.  The reason\hand{|}is surely an inner reluctance,
an unreadiness to embark upon a long-term voyage,  well knowing that it is
\inserted{so} enticing \xout{xxxxxxx,}that I may well be caught in this game for a number of
years - not doing anything else \xout{from} day and night than making love with
mathematics,  and maybe sleeping and eating now and then.  I have gone
through this a number of times, \inserted{and} at times I thought I was through.  Finally
I came to admit and to accept,  two years ago,  I was not through yet -
this was during the months of meditation after \unsure{that|the} ``long march through
Galois theory'' -  which had been\hand{,} too\hand{,} a wholly unexpected fit of mathemati-%
cal passion, not to say frenzy. \deleted{Finally} \inserted{And} \unsure{\over{D}{d}}uring the last weeks,  just
reflecting a little here and there upon the Teichmüller-Légo game and
it's arithmetical implications,  I let myself be caught again by this
fascination \xout{ofxxx} - it is becoming kind of clear \inserted{now} that I am going to finish
writing up those notes on algebra,  almost like some homework that has got
to be done (anyhow I like to finish when I started \xout{xx} something)  -  and
as soon as I'm through with the notes,  back to geom\unsure{\over{w}{e}}try in the long last !
Also,  the idea is in the air for the last few months  -  since I decided
to publish these unformal notes on stacks or whatever it'll turn out to be -
that I may well \inserted{go on} the same way,  writing up and publishing unformal notes
on other topics,  including tame topology and anabelian algebraic geometry.
In contrast to the present notes,  I got heaps of scratchwork done \inserted{on these} in the
years before  -  in this respect time is even riper for me to ramble \inserted{``publicly'',} \deleted{on}
\deleted{\xout{these};} than on stacks and homotopy theory\unsure{\over{.}{!}}

From Y.  who looked through a lot of litterature on the subject,  it
strikes me (agreeably of course) that nobody yet hit upon ``the'' natural
presentation of the Teichmüller groupoids,  which kind of imposes itself
quite \over{naturall}{forcibl}y in the set-up I let myself be guided by.  Technically spea-%
king (and this will rejoice Ronnie Brown I'm sure!),  I suspect one main rea-%
son why this is so,  is that people are accustomed to work with fundamental
groups and generators and relations for these and stick to it,  even in
contexts when this is wholly unadequate,  namely when you get a clear

\margin{7}
